\chapter{Dualità}\label{ch:dualita}
La teoria della \textbf{dualità} associa a ogni problema di PL, detto \textbf{primale}, un altro problema di PL, detto \textbf{duale}. Il duale si costruisce con gli stessi dati del primale (la matrice $A$ e i vettori $b$ e $c$), ma ha variabili diverse e un obiettivo opposto: a un primale di minimo corrisponde un duale di massimo, e viceversa. Il concetto di dualità è stato introdotto da Von Neumann nel 1947, mentre il teorema della dualità è stato dimostrato in modo rigoroso nel 1951 da Gale, Kuhn e Tucker.

Si potrebbe pensare che affiancare un secondo problema al primo serva solo a complicare le cose. In realtà lo studio del duale permette di:
\begin{itemize}
    \item \textbf{Mettere in luce proprietà del primale}: ad esempio il duale dà un valore economico alle risorse del problema di partenza (sezione~\ref{sec:duale-economico}).
    \item \textbf{Risolvere un problema più facile}: a volte il duale è più semplice da risolvere del primale, e risolvendo l'uno si ottengono informazioni sull'altro.
    \item \textbf{Costruire un'altra versione del simplesso}: il \textbf{simplesso duale}, utile ad esempio quando la soluzione di base ha componenti negative (sezione~\ref{sec:simplesso-duale}).
\end{itemize}

\section{Il problema duale}
\begin{definizione}[Problema duale]\label{def:duale}
Dato il problema primale
\[
\min c^T x, \quad A x \ge b, \quad x \ge 0 \tag{P}
\]
con $A \in \mathbb{R}^{m \times n}$, $b \in \mathbb{R}^m$ e $c \in \mathbb{R}^n$, si dice \textbf{problema duale} il problema
\[
\max y^T b, \quad y^T A \le c^T, \quad y \ge 0 \tag{D}
\]
dove $y \in \mathbb{R}^m$ è il vettore delle \textbf{variabili duali}. I due problemi formano una coppia di \textbf{problemi duali simmetrici}.
\end{definizione}

Il vincolo $y^T A \le c^T$ si legge colonna per colonna: la $j$-esima componente di $y^T A$ è il prodotto scalare tra $y$ e la $j$-esima colonna di $A$. Per esteso, il duale è quindi:
\[
\begin{aligned}
&\max \sum_{i=1}^{m} b_i y_i \\
&\sum_{i=1}^{m} a_{ij} y_i \le c_j, \quad j = 1, \dots, n \\
&y_i \ge 0, \quad i = 1, \dots, m.
\end{aligned}
\]
Confrontando i due problemi si nota che:
\begin{itemize}
    \item \textbf{Variabili e vincoli si scambiano}: il primale ha $n$ variabili e $m$ vincoli, il duale ha $m$ variabili e $n$ vincoli. A ogni vincolo $i$ del primale corrisponde una variabile duale $y_i$, e a ogni variabile $x_j$ del primale corrisponde un vincolo duale.
    \item \textbf{Costi e termini noti si scambiano}: i coefficienti $c_j$ della funzione obiettivo primale diventano i termini noti del duale, e i termini noti $b_i$ del primale diventano i coefficienti della funzione obiettivo duale.
    \item \textbf{La matrice si traspone}: il $j$-esimo vincolo duale usa la $j$-esima colonna di $A$, quindi la matrice dei vincoli duali è $A^T$.
\end{itemize}

\paragraph{Esempio.}
Consideriamo il problema primale
\[
\begin{aligned}
&\min\, (2x_1 + x_2) \\
&2x_1 + 3x_2 + x_3 \ge 4 \\
&x_1 + 5x_2 + 4x_3 \ge 8 \\
&x_1,\, x_2,\, x_3 \ge 0
\end{aligned}
\qquad
A = \begin{pmatrix} 2 & 3 & 1 \\ 1 & 5 & 4 \end{pmatrix}, \quad
b = \begin{pmatrix} 4 \\ 8 \end{pmatrix}, \quad
c = \begin{pmatrix} 2 \\ 1 \\ 0 \end{pmatrix}
\]
con $n = 3$ variabili e $m = 2$ vincoli. La variabile $x_3$ non compare nella funzione obiettivo, quindi il suo costo è $c_3 = 0$. Il duale ha due variabili $y_1, y_2$, una per ogni vincolo, e tre vincoli, uno per ogni colonna di $A$:
\[
\begin{aligned}
&\max\, (4y_1 + 8y_2) \\
&2y_1 + y_2 \le 2 \\
&3y_1 + 5y_2 \le 1 \\
&y_1 + 4y_2 \le 0 \\
&y_1,\, y_2 \ge 0
\end{aligned}
\qquad
A^T = \begin{pmatrix} 2 & 1 \\ 3 & 5 \\ 1 & 4 \end{pmatrix}.
\]
Il primo vincolo duale usa la colonna di $x_1$, ovvero $(2, 1)^T$, e ha come termine noto il costo $c_1 = 2$. Allo stesso modo il secondo usa la colonna di $x_2$, con termine noto $c_2 = 1$, e il terzo quella di $x_3$, con termine noto $c_3 = 0$.

\begin{osservazione}[Il duale del duale]
Il duale del problema duale è il problema primale. Per vederlo, scriviamo (D) come problema di minimo con vincoli di tipo $\ge$:
\[
\min\, (-b)^T y, \quad (-A^T)\, y \ge -c, \quad y \ge 0.
\]
Applicando la definizione, il suo duale è $\max\, (-c)^T x$ con $x^T(-A^T) \le -b^T$ e $x \ge 0$, ovvero $\min c^T x$ con $A x \ge b$ e $x \ge 0$: è proprio (P). Primale e duale sono quindi uno il duale dell'altro, e non importa quale dei due chiamiamo primale.
\end{osservazione}

\section{Il duale della forma standard}\label{sec:duale-standard}
Se il primale è in forma standard, i vincoli sono di uguaglianza:
\[
\min c^T x, \quad A x = b, \quad x \ge 0. \tag{P}
\]
In questo caso il duale è
\[
\max y^T b, \quad y^T A \le c^T \tag{D}
\]
dove le variabili duali $y$ non hanno vincoli di segno, ovvero sono \textbf{libere}. Questa coppia si dice di \textbf{problemi duali asimmetrici}.

Il motivo è che un vincolo di uguaglianza equivale a due disuguaglianze: $Ax = b$ significa $Ax \ge b$ e $-Ax \ge -b$. Il primale è quindi nella forma simmetrica, con matrice $\begin{pmatrix} A \\ -A \end{pmatrix}$ e termini noti $\begin{pmatrix} b \\ -b \end{pmatrix}$. Applicando la definizione con due vettori di variabili duali $u, v \ge 0$, uno per ogni gruppo di vincoli, si ottiene:
\[
\max\, (u^T b - v^T b), \quad u^T A - v^T A \le c^T, \quad u, v \ge 0.
\]
Ponendo $y = u - v$, che può avere componenti di qualsiasi segno, si ritrova proprio (D).

\section{Regole di trasformazione}\label{sec:regole-duale}
In generale un problema di PL può avere vincoli di tutti i tipi e variabili con segni diversi. Per scriverne il duale non serve portarlo prima nella forma (P): si usano direttamente le regole della tabella~\ref{tab:trasformazioni}. La tabella si legge da sinistra a destra quando il primale è di minimo, e da destra a sinistra quando il primale è di massimo.

\begin{table}[htbp]
\centering
\begin{tabular}{ll}
\toprule
\textbf{Problema di minimo} & \textbf{Problema di massimo} \\
\midrule
coefficienti della funzione obiettivo & termini noti \\
termini noti & coefficienti della funzione obiettivo \\
coefficienti dell'$i$-esimo vincolo & coefficienti della $i$-esima variabile \\
\midrule
$i$-esimo vincolo $\ge$ & $i$-esima variabile $\ge 0$ \\
$i$-esimo vincolo $\le$ & $i$-esima variabile $\le 0$ \\
$i$-esimo vincolo $=$ & $i$-esima variabile libera \\
\midrule
$j$-esima variabile $\ge 0$ & $j$-esimo vincolo $\le$ \\
$j$-esima variabile $\le 0$ & $j$-esimo vincolo $\ge$ \\
$j$-esima variabile libera & $j$-esimo vincolo $=$ \\
\midrule
numero di vincoli & numero di variabili \\
numero di variabili & numero di vincoli \\
\bottomrule
\end{tabular}
\caption{Regole di trasformazione tra un problema di minimo e il suo duale, che è un problema di massimo. Ogni riga mette in corrispondenza un elemento del primo problema con l'elemento corrispondente del secondo.}
\label{tab:trasformazioni}
\end{table}

Le regole sono più facili da ricordare se si pensa alla coppia simmetrica (P) e (D). In un problema di minimo la forma \emph{naturale} ha vincoli $\ge$ e variabili $\ge 0$, mentre in un problema di massimo ha vincoli $\le$ e variabili $\ge 0$. Un elemento in forma naturale corrisponde a un elemento in forma naturale, uno \emph{non naturale} (come un vincolo $\le$ in un problema di minimo) corrisponde a uno non naturale, e un'uguaglianza corrisponde a una variabile libera.

In pratica, il tipo di ogni vincolo primale determina il segno della variabile duale corrispondente, mentre il segno di ogni variabile primale determina il tipo del vincolo duale corrispondente.

\paragraph{Esempio.}
Scriviamo il duale del problema di massimo
\[
\begin{aligned}
&\max\, (24x_1 + 14x_2) \\
&3x_1 + 2x_2 \le 12 \\
&4x_1 + x_2 \le 10 \\
&2x_1 + x_2 \le 7 \\
&x_1,\, x_2 \ge 0
\end{aligned}
\qquad
A = \begin{pmatrix} 3 & 2 \\ 4 & 1 \\ 2 & 1 \end{pmatrix}.
\]
Il primale è di massimo, quindi leggiamo la tabella da destra a sinistra. I tre vincoli sono di tipo $\le$, quindi le tre variabili duali sono $\ge 0$. Le due variabili primali sono $\ge 0$, quindi i due vincoli duali sono di tipo $\ge$:
\[
\begin{aligned}
&\min\, (12y_1 + 10y_2 + 7y_3) \\
&3y_1 + 4y_2 + 2y_3 \ge 24 \\
&2y_1 + y_2 + y_3 \ge 14 \\
&y_1,\, y_2,\, y_3 \ge 0
\end{aligned}
\qquad
A^T = \begin{pmatrix} 3 & 4 & 2 \\ 2 & 1 & 1 \end{pmatrix}.
\]
Ritroveremo questo problema nella sezione~\ref{sec:post-ottimale}, con i termini noti moltiplicati per $100$.

\paragraph{Esempio.}
Scriviamo il duale del problema
\[
\begin{aligned}
&\min\, (2x_1 + 3x_2 + 4x_3 + x_4) \\
&x_1 - 5x_3 + 2x_4 \ge 7 \\
&2x_1 + 4x_2 - 6x_3 \ge 9 \\
&x_1,\, x_2,\, x_3,\, x_4 \text{ libere}
\end{aligned}
\qquad
A = \begin{pmatrix} 1 & 0 & -5 & 2 \\ 2 & 4 & -6 & 0 \end{pmatrix}.
\]
I due vincoli sono di tipo $\ge$, quindi $y_1, y_2 \ge 0$. Le variabili primali invece sono tutte libere, quindi i quattro vincoli duali sono di uguaglianza:
\[
\begin{aligned}
&\max\, (7y_1 + 9y_2) \\
&y_1 + 2y_2 = 2 \\
&4y_2 = 3 \\
&{-5y_1} - 6y_2 = 4 \\
&2y_1 = 1 \\
&y_1,\, y_2 \ge 0.
\end{aligned}
\]
Torneremo su questo duale nella sezione~\ref{sec:teoremi-dualita}.

\paragraph{Esempio.}
Scriviamo il duale del problema
\[
\begin{aligned}
&\min\, (2x_1 + 3x_2 + x_3) \\
&3x_1 + x_2 + 5x_3 \ge 7 \\
&x_1 + x_2 - 6x_3 \le 9 \\
&4x_1 - x_2 - 2x_3 = 8 \\
&x_1,\, x_2,\, x_3 \ge 0
\end{aligned}
\qquad
A = \begin{pmatrix} 3 & 1 & 5 \\ 1 & 1 & -6 \\ 4 & -1 & -2 \end{pmatrix}.
\]
Il primale è di minimo, quindi leggiamo la tabella da sinistra a destra. Le variabili duali sono tre, una per vincolo, e il loro segno dipende dal tipo di vincolo:
\begin{itemize}
    \item Il primo vincolo è $\ge$, quindi $y_1 \ge 0$.
    \item Il secondo vincolo è $\le$, quindi $y_2 \le 0$.
    \item Il terzo vincolo è di uguaglianza, quindi $y_3$ è libera.
\end{itemize}
Le variabili primali sono invece tutte $\ge 0$, quindi i tre vincoli duali, uno per ogni colonna di $A$, sono di tipo $\le$:
\[
\begin{aligned}
&\max\, (7y_1 + 9y_2 + 8y_3) \\
&3y_1 + y_2 + 4y_3 \le 2 \\
&y_1 + y_2 - y_3 \le 3 \\
&5y_1 - 6y_2 - 2y_3 \le 1 \\
&y_1 \ge 0, \quad y_2 \le 0, \quad y_3 \text{ libera.}
\end{aligned}
\]

\section{Interpretazione economica}\label{sec:duale-economico}
Le variabili duali hanno spesso un significato concreto. Riprendiamo il problema di pianificazione della produzione del capitolo~\ref{ch:pl}, in cui si vuole massimizzare il profitto realizzando $n$ prodotti con $m$ risorse limitate:
\[
\max \sum_{j=1}^{n} c_j x_j, \qquad \sum_{j=1}^{n} a_{ij} x_j \le b_i, \quad i = 1, \dots, m, \qquad x_j \ge 0.
\]
Il primale è di massimo, quindi leggiamo la tabella~\ref{tab:trasformazioni} da destra a sinistra: i vincoli $\le$ danno variabili duali $y_i \ge 0$, e le variabili $x_j \ge 0$ danno vincoli duali di tipo $\ge$. Il duale è quindi:
\[
\min \sum_{i=1}^{m} b_i y_i, \qquad \sum_{i=1}^{m} a_{ij} y_i \ge c_j, \quad j = 1, \dots, n, \qquad y_i \ge 0.
\]

Supponiamo ora che, in un mercato parallelo, qualcuno voglia comprare tutte le nostre risorse invece di lasciarci produrre, e indichiamo con $y_i$ il prezzo unitario che offre per la risorsa $R_i$. Comprando tutte le risorse spende $\sum_i b_i y_i$, e per essere competitivo vuole spendere il meno possibile. Per convincerci a vendere, però, deve offrire prezzi convenienti: per ogni prodotto $P_j$, le risorse necessarie a fabbricarne una unità devono essere pagate almeno quanto il profitto $c_j$ che ne ricaveremmo, ovvero $\sum_i a_{ij} y_i \ge c_j$.

Il problema del compratore è proprio il duale. Anche il duale è quindi un problema di natura economica, e il mercato parallelo si dice \textbf{mercato ombra}. I valori $y_i$ si chiamano \textbf{prezzi ombra} delle risorse, perché misurano il valore implicito di ciascuna risorsa all'interno del problema.

\paragraph{Esempio.}
Torniamo al problema dei mattoncini lego~\eqref{eq:lego}:
\[
\begin{aligned}
&\max\, (15x_1 + 20x_2) \\
&2x_1 + 2x_2 \le 8 \\
&x_1 + 2x_2 \le 6 \\
&x_1,\, x_2 \ge 0
\end{aligned}
\qquad\Longrightarrow\qquad
\begin{aligned}
&\min\, (8y_1 + 6y_2) \\
&2y_1 + y_2 \ge 15 \\
&2y_1 + 2y_2 \ge 20 \\
&y_1,\, y_2 \ge 0
\end{aligned}
\]
dove $y_1$ e $y_2$ sono i prezzi offerti per un mattoncino piccolo e per uno grande. Il primo vincolo duale dice che i mattoncini di una sedia (2 piccoli e 1 grande) devono valere almeno quanto la sedia, ovvero 15 euro. Il secondo dice lo stesso per il tavolo (2 piccoli e 2 grandi, 20 euro).

Il duale ha due variabili, quindi possiamo risolverlo graficamente (figura~\ref{fig:duale-lego}). La regione ammissibile non è limitata e ha tre vertici: $(0, 15)$, $(5, 5)$ e $(10, 0)$, dove il vertice $(5, 5)$ è l'intersezione delle rette $2y_1 + y_2 = 15$ e $y_1 + y_2 = 10$. La funzione obiettivo vale rispettivamente $90$, $70$ e $80$, e nelle direzioni in cui la regione si estende cresce. Il minimo è quindi $70$, con $y_1 = y_2 = 5$ euro a mattoncino.

Notiamo che $70$ euro è proprio il profitto massimo del primale, che si ottiene costruendo 2 sedie e 2 tavoli (sezione~\ref{sec:graf}). Non è un caso: il fatto che i valori ottimi del primale e del duale coincidano è il contenuto del teorema della dualità forte, che vediamo nella prossima sezione.

\begin{figure}[htbp]
\centering
\begin{tikzpicture}[scale=0.32, >=Stealth]
  \useasboundingbox (-11,-2.6) rectangle (21,20.0);
  % regione ammissibile (non limitata)
  \fill[Accent!12] (0,17) -- (0,15) -- (5,5) -- (10,0) -- (17,0) -- (17,17) -- cycle;
  % assi
  \draw[->] (-1,0) -- (18.5,0) node[right] {$y_1$};
  \draw[->] (0,-1) -- (0,18.5) node[above] {$y_2$};
  \foreach \t in {5,10,15} {
    \draw (\t,0.3) -- (\t,-0.3);
    \draw (0.3,\t) -- (-0.3,\t);
  }
  % vincoli
  \draw[blue!65!black, very thick] (-0.5,16) -- (8,-1)
    node[pos=0.22, sloped, above, font=\scriptsize] {$2y_1+y_2=15$};
  \draw[orange!85!black, very thick] (-0.5,10.5) -- (11,-1)
    node[pos=0.7, sloped, above, font=\scriptsize] {$y_1+y_2=10$};
  % curva di livello ottima
  \draw[dashed, red!70!black] (-0.5,12.33) -- (9.5,-1);
  \node[red!70!black, font=\scriptsize, anchor=east] at (-0.7,12.33) {$8y_1+6y_2=70$};
  % direzione di decrescita
  \draw[->, red!70!black, very thick] (14,10) -- (11.6,8.2) node[below left, font=\scriptsize] {$d$};
  % vertici
  \node[circle, fill=black, inner sep=1.4pt] at (0,15) {};
  \node[circle, fill=red!70!black, inner sep=1.8pt] at (5,5) {};
  \node[circle, fill=black, inner sep=1.4pt] at (10,0) {};
  \node[font=\scriptsize, anchor=east] at (-0.5,15) {$(0,15)$};
  \node[font=\scriptsize, anchor=north east, red!70!black] at (4.6,4.6) {$(5,5)$};
  \node[font=\scriptsize, anchor=north] at (10,-0.5) {$(10,0)$};
\end{tikzpicture}
\caption{Risoluzione grafica del duale del problema dei mattoncini: la regione ammissibile (zona colorata) non è limitata ed è delimitata dalle rette dei due vincoli e dagli assi, con vertici $(0,15)$, $(5,5)$ e $(10,0)$. Spostando le curve di livello $8y_1+6y_2=k$ nella direzione di decrescita $d=(-8,-6)^T$, l'ultimo punto della regione toccato è il vertice $(5,5)$, sulla curva tratteggiata $8y_1+6y_2=70$.}
\label{fig:duale-lego}
\end{figure}

\section{Teoremi della dualità}\label{sec:teoremi-dualita}
In questa sezione consideriamo la coppia di problemi duali simmetrici
\[
\min c^T x, \quad A x \ge b, \quad x \ge 0 \tag{P}
\]
\[
\max y^T b, \quad y^T A \le c^T, \quad y \ge 0. \tag{D}
\]
Il primo risultato confronta i valori delle due funzioni obiettivo in punti ammissibili qualsiasi.

\begin{teorema}[Dualità debole]\label{teo:dualita-debole}
Se $\bar{x}$ è un punto ammissibile per il primale e $\bar{y}$ è un punto ammissibile per il duale, allora
\[
\bar{y}^T b \le c^T \bar{x}.
\]
\end{teorema}

\paragraph{Dimostrazione.}
Poiché $\bar{x}$ è ammissibile si ha $A \bar{x} \ge b$, e poiché $\bar{y} \ge 0$ possiamo moltiplicare la disuguaglianza per $\bar{y}^T$ senza cambiarne il verso: $\bar{y}^T b \le \bar{y}^T A \bar{x}$. Allo stesso modo $\bar{y}$ è ammissibile, quindi $\bar{y}^T A \le c^T$, e poiché $\bar{x} \ge 0$ si ha $\bar{y}^T A \bar{x} \le c^T \bar{x}$. Mettendo insieme le due disuguaglianze:
\begin{equation}\label{eq:catena-debole}
\bar{y}^T b \le \bar{y}^T A \bar{x} \le c^T \bar{x}
\end{equation}
che è la tesi. \hfill $\square$

\vspace{0.5em}
In pratica, il duale in un suo punto ammissibile non supera mai il valore del primale in un suo punto ammissibile. Ne seguono alcune conseguenze:
\begin{itemize}
    \item Il valore della funzione obiettivo primale in un qualsiasi punto ammissibile è un \textbf{upper bound} (limite superiore) per il valore ottimo del duale.
    \item Il valore della funzione obiettivo duale in un qualsiasi punto ammissibile è un \textbf{lower bound} (limite inferiore) per il valore ottimo del primale.
    \item Se primale e duale sono entrambi ammissibili, i loro valori ottimi sono finiti, e l'ottimo del primale è maggiore o uguale dell'ottimo del duale.
\end{itemize}

\begin{corollario}[Condizione di ottimalità]\label{cor:ottimalita-duale}
Se $\bar{x}$ è ammissibile per il primale, $\bar{y}$ è ammissibile per il duale e $c^T \bar{x} = \bar{y}^T b$, allora $\bar{x}$ è una soluzione ottima del primale e $\bar{y}$ è una soluzione ottima del duale.
\end{corollario}

Il motivo è che, per la dualità debole, nessun punto ammissibile del primale può avere un valore minore di $\bar{y}^T b$, e $\bar{x}$ raggiunge proprio questo valore. Lo stesso ragionamento vale per $\bar{y}$. Il corollario dà quindi un modo per certificare che una soluzione è ottima: basta trovare una soluzione duale ammissibile con lo stesso valore.

\begin{corollario}\label{cor:illimitato-inammissibile}
Se il primale è illimitato inferiormente, allora il duale è inammissibile. Allo stesso modo, se il duale è illimitato superiormente, allora il primale è inammissibile.
\end{corollario}

Infatti, se il duale avesse un punto ammissibile $\bar{y}$, il valore $\bar{y}^T b$ sarebbe un lower bound per il primale, che quindi non potrebbe essere illimitato inferiormente.

La dualità debole dice che l'ottimo del duale non supera quello del primale. Il teorema della dualità forte dice che, in realtà, i due valori coincidono.

\begin{teorema}[Dualità forte]\label{teo:dualita-forte}
Se uno dei due problemi, primale o duale, ammette una soluzione ottima finita, anche l'altro ammette una soluzione ottima finita, e i valori ottimi delle due funzioni obiettivo coincidono. Se invece uno dei due problemi è illimitato, l'altro non ammette soluzioni ammissibili\footnote{In questa forma, sia la dualità debole sia la dualità forte valgono per i problemi di PL.}.
\end{teorema}

Mettendo insieme i due teoremi si ottengono tutte le combinazioni possibili tra primale e duale, riassunte nella tabella~\ref{tab:primale-duale}. Ad esempio, un primale con ottimo finito non può avere un duale illimitato o inammissibile, mentre primale e duale possono essere entrambi inammissibili.

\begin{table}[htbp]
\centering
\begin{tabular}{lccc}
\toprule
 & \multicolumn{3}{c}{\textbf{Duale}} \\
\cmidrule(lr){2-4}
\textbf{Primale} & ha ottimo finito & è illimitato & è inammissibile \\
\midrule
ha ottimo finito & sì & no & no \\
è illimitato & no & no & sì \\
è inammissibile & no & sì & sì \\
\bottomrule
\end{tabular}
\caption{Relazioni tra primale e duale: ogni casella dice se il primale (riga) e il duale (colonna) possono trovarsi contemporaneamente nei casi indicati.}
\label{tab:primale-duale}
\end{table}

\paragraph{Esempio.}
Torniamo al duale del problema con variabili libere della sezione~\ref{sec:regole-duale}. Le ultime due equazioni danno $y_1 = \frac12$ e $y_2 = \frac34$, che soddisfano anche la prima, perché $\frac12 + \frac32 = 2$. Con questi valori, però, la terza equazione vale $-\frac52 - \frac92 = -7 \neq 4$: il duale non ha soluzioni ammissibili. Il primale invece è ammissibile, ad esempio con $x = (7, 0, 0, 0)$. Guardando la tabella~\ref{tab:primale-duale}, un duale inammissibile è compatibile solo con un primale illimitato o inammissibile: il primale è quindi illimitato inferiormente, senza bisogno di risolverlo.

Per un primale in forma standard, la soluzione ottima duale si può costruire direttamente a partire dalla soluzione ottima di base del primale, come dice il teorema seguente.

\begin{teorema}[Soluzione ottima duale]\label{teo:sol-duale}
Sia dato il primale in forma standard $\min c^T x$, $Ax = b$, $x \ge 0$, e sia $\bar{x} = (B^{-1}b, 0)$ una sua soluzione ottima di base, con costi ridotti non negativi. Allora
\[
\bar{y}^T = c_B^T B^{-1}
\]
è una soluzione ottima del duale $\max y^T b$, $y^T A \le c^T$.
\end{teorema}

\paragraph{Dimostrazione.}
Verifichiamo prima che $\bar{y}$ sia ammissibile per il duale, ovvero che $\bar{y}^T A \le c^T$. Partizionando $A = (B, N)$, il vincolo si divide in due parti. Per le colonne di base si ha
\[
\bar{y}^T B = c_B^T B^{-1} B = c_B^T
\]
quindi il vincolo è soddisfatto con l'uguaglianza. Per le colonne fuori base si ha invece
\[
\bar{y}^T N = c_B^T B^{-1} N \le c_N^T
\]
perché $c_N^T - c_B^T B^{-1} N \ge 0$ è proprio la condizione sui costi ridotti. Inoltre
\[
\bar{y}^T b = c_B^T B^{-1} b = c^T \bar{x}
\]
quindi primale e duale hanno lo stesso valore, e per il corollario~\ref{cor:ottimalita-duale}\footnote{Il corollario vale anche per la coppia asimmetrica, perché la dualità debole si dimostra allo stesso modo: se $A\bar{x} = b$, $\bar{x} \ge 0$ e $\bar{y}^T A \le c^T$, allora $\bar{y}^T b = \bar{y}^T A \bar{x} \le c^T \bar{x}$.} $\bar{y}$ è ottima. \hfill $\square$

\begin{osservazione}[Ammissibilità duale e ottimalità primale]
La dimostrazione mostra che $\bar{y}^T = c_B^T B^{-1}$ è ammissibile per il duale se e solo se i costi ridotti della base $B$ sono non negativi. In altre parole, la condizione di ammissibilità della soluzione duale coincide con la condizione di ottimalità primale (teorema~\ref{teo:ottimalita}). È l'idea su cui si basa il simplesso duale.
\end{osservazione}

\section{Condizioni di complementarietà}\label{sec:complementarieta}
La dualità forte dice che all'ottimo i valori di primale e duale coincidono. Le \textbf{condizioni di complementarietà}, dette anche degli \textbf{scarti complementari}, traducono questa uguaglianza in condizioni sui singoli vincoli.

\begin{teorema}[Scarti complementari]\label{teo:complementarieta}
Siano $\bar{x}$ ammissibile per (P) e $\bar{y}$ ammissibile per (D). Allora $\bar{x}$ e $\bar{y}$ sono soluzioni ottime dei rispettivi problemi se e solo se
\[
\bar{y}^T (A\bar{x} - b) = 0 \qquad \text{e} \qquad (\bar{y}^T A - c^T)\, \bar{x} = 0.
\]
\end{teorema}

Il motivo segue dalla catena di disuguaglianze~\eqref{eq:catena-debole}. La prima disuguaglianza è un'uguaglianza solo se $\bar{y}^T(A\bar{x} - b) = 0$, mentre la seconda lo è solo se $(\bar{y}^T A - c^T)\bar{x} = 0$. I due punti sono ottimi se e solo se $\bar{y}^T b = c^T \bar{x}$ (per la dualità forte e il corollario~\ref{cor:ottimalita-duale}), ovvero se e solo se valgono entrambe le uguaglianze.

Le due condizioni si possono scrivere componente per componente. Il vettore $A\bar{x} - b$ ha componenti $\ge 0$, perché $\bar{x}$ è ammissibile, e anche $\bar{y} \ge 0$: la somma $\bar{y}^T(A\bar{x} - b)$ è quindi nulla solo se ogni suo termine è nullo. Lo stesso vale per la seconda condizione, in cui $\bar{y}^T A - c^T \le 0$ e $\bar{x} \ge 0$. Si ottengono così $m + n$ condizioni:
\[
\bar{y}_i \,(a_i^T \bar{x} - b_i) = 0, \quad i = 1, \dots, m \qquad\qquad \bar{x}_j \,(\bar{y}^T A_j - c_j) = 0, \quad j = 1, \dots, n
\]
dove $a_i^T$ è la $i$-esima riga di $A$ e $A_j$ è la sua $j$-esima colonna.

In pratica, per ogni vincolo almeno uno dei due fattori deve essere nullo: se $\bar{y}_i > 0$ il vincolo $i$ è attivo, e se il vincolo $i$ non è attivo allora $\bar{y}_i = 0$. Lo stesso vale per ogni variabile primale e il vincolo duale corrispondente.

Pensando ai prezzi ombra della sezione~\ref{sec:duale-economico}, la condizione ha una lettura molto banale. Se il prezzo ombra di una risorsa è positivo, la risorsa è tutta utilizzata (il vincolo è saturo). Se invece una parte della risorsa avanza, il suo prezzo ombra è nullo: averne di più non serve.

\begin{osservazione}[Casi particolari]
Se il primale è in forma standard, i vincoli $Ax = b$ sono di uguaglianza e la prima condizione è sempre verificata. Le condizioni di complementarietà si riducono quindi a $(\bar{y}^T A - c^T)\bar{x} = 0$. Al contrario, se il primale ha vincoli $Ax \ge b$ e variabili libere, il duale ha vincoli di uguaglianza $y^T A = c^T$, e resta solo la condizione $\bar{y}^T(A\bar{x} - b) = 0$.
\end{osservazione}

\paragraph{Esempio.}
Consideriamo il problema
\[
\begin{aligned}
&\min\, (6x_1 + 4x_2) \\
&x_1 + 2x_2 \ge 1 \\
&2x_1 + x_2 \ge -1 \\
&x_1,\, x_2 \ge 0
\end{aligned}
\]
di cui sappiamo che la soluzione ottima è $\bar{x} = (0, \frac12)$, e calcoliamo la soluzione ottima duale. Per prima cosa costruiamo il duale:
\[
\begin{aligned}
&\max\, (y_1 - y_2) \\
&y_1 + 2y_2 \le 6 \\
&2y_1 + y_2 \le 4 \\
&y_1,\, y_2 \ge 0.
\end{aligned}
\]
Le condizioni di complementarietà sono quattro, due per i vincoli primali e due per le variabili primali:
\[
\begin{aligned}
y_1\,(x_1 + 2x_2 - 1) &= 0 \\
y_2\,(2x_1 + x_2 + 1) &= 0
\end{aligned}
\qquad\qquad
\begin{aligned}
x_1\,(y_1 + 2y_2 - 6) &= 0 \\
x_2\,(2y_1 + y_2 - 4) &= 0.
\end{aligned}
\]
Sostituiamo $\bar{x} = (0, \frac12)$. Il primo vincolo primale è attivo, quindi la prima condizione è soddisfatta qualunque sia $y_1$. Il secondo vale invece $\frac12 + 1 = \frac32 > 0$, quindi deve essere $y_2 = 0$. Poiché $x_1 = 0$, la terza condizione è sempre soddisfatta, mentre $x_2 = \frac12 > 0$ impone $2y_1 + y_2 = 4$, da cui $y_1 = 2$. La soluzione ottima duale è quindi $\bar{y} = (2, 0)$. Come verifica, i due valori ottimi coincidono: $6 \cdot 0 + 4 \cdot \frac12 = 2 = 2 - 0$.

\paragraph{Esempio.}
Torniamo al problema dei mattoncini e al suo duale. Il primale è di massimo, ma le condizioni hanno la stessa forma: ogni variabile duale moltiplica lo scarto del vincolo primale corrispondente, e ogni variabile primale moltiplica lo scarto del vincolo duale corrispondente.

La soluzione ottima primale è $x^* = (2, 2)$, quindi entrambe le variabili sono positive e i due vincoli duali devono essere attivi:
\[
\begin{cases}
2y_1 + y_2 = 15 \\
2y_1 + 2y_2 = 20
\end{cases}
\quad\Longrightarrow\quad y_1 = y_2 = 5.
\]
Ritroviamo così, senza disegnare nulla, la soluzione trovata graficamente. Le condizioni sui vincoli primali, invece, non danno informazioni: nella soluzione ottima si usano tutti i mattoncini, quindi entrambi i vincoli sono attivi e le condizioni sono soddisfatte qualunque siano $y_1$ e $y_2$.

\section{Calcolo della soluzione duale dalla tabella}\label{sec:duale-tabella}
Il teorema~\ref{teo:sol-duale} dà anche un metodo operativo per calcolare la soluzione ottima duale a partire dalla tabella finale del simplesso applicato al primale:
\[
\bar{y}^T = c_B^T B^{-1}
\]
dove $c_B$ è il vettore dei costi delle variabili in base nella tabella finale, ordinati come le righe della tabella.

Per trovare $B^{-1}$ non serve invertire nessuna matrice. Le operazioni del simplesso trasformano infatti i vincoli della tabella iniziale moltiplicandoli per $B^{-1}$, quindi le colonne che nella tabella iniziale formavano la matrice identità diventano $B^{-1} I = B^{-1}$. In pratica, $B^{-1}$ si legge nella tabella finale, nelle colonne delle variabili che erano in base alla prima iterazione.

Se la base iniziale è formata da variabili di scarto si può fare ancora più in fretta. Il costo ridotto della variabile di scarto dell'$i$-esimo vincolo, che ha costo nullo e colonna $e_i$, vale infatti $0 - \bar{y}^T e_i = -\bar{y}_i$. Quindi, se il primale ha solo vincoli di tipo $\le$:
\[
\bar{y}_i = -\bar{c}_{\text{scarto } i}
\]
ovvero la soluzione duale è l'opposto del vettore dei costi ridotti delle variabili di scarto nella tabella finale.

\paragraph{Esempio.}
Consideriamo il problema
\[
\begin{aligned}
&\min\, (-2x_1 - x_2) \\
&x_1 + 2x_2 \le 2 \\
&2x_1 + 3x_2 \le 4 \\
&x_1 + 3x_2 \le 3 \\
&x_1,\, x_2 \ge 0
\end{aligned}
\qquad\Longrightarrow\qquad
\begin{aligned}
&\max\, (2y_1 + 4y_2 + 3y_3) \\
&y_1 + 2y_2 + y_3 \le -2 \\
&2y_1 + 3y_2 + 3y_3 \le -1 \\
&y_1,\, y_2,\, y_3 \le 0
\end{aligned}
\]
dove il duale ha variabili $\le 0$ perché i vincoli del primale di minimo sono di tipo $\le$. Risolviamo il primale con il simplesso, aggiungendo le variabili di scarto $x_3$, $x_4$ e $x_5$. Entra $x_1$, e i rapporti delle righe di $x_3$ e di $x_4$ valgono entrambi $2$: per la regola dell'indice minimo esce $x_3$. Con le operazioni $R_2 \to R_2 - 2R_1$, $R_3 \to R_3 - R_1$ e $R_4 \to R_4 + 2R_1$ si ottiene:
\begin{center}
\begin{tabular}{c|ccccc|c}
 & $x_1$ & $x_2$ & $x_3$ & $x_4$ & $x_5$ & $b$ \\
\hline
$x_3$ & $\boxed{1}$ & 2 & 1 & 0 & 0 & 2 \\
$x_4$ & 2 & 3 & 0 & 1 & 0 & 4 \\
$x_5$ & 1 & 3 & 0 & 0 & 1 & 3 \\
\hline
 & $-2$ & $-1$ & 0 & 0 & 0 & 0 \\
\end{tabular}
\qquad$\Longrightarrow$\qquad
\begin{tabular}{c|ccccc|c}
 & $x_1$ & $x_2$ & $x_3$ & $x_4$ & $x_5$ & $b$ \\
\hline
$x_1$ & 1 & 2 & 1 & 0 & 0 & 2 \\
$x_4$ & 0 & $-1$ & $-2$ & 1 & 0 & 0 \\
$x_5$ & 0 & 1 & $-1$ & 0 & 1 & 1 \\
\hline
 & 0 & 3 & 2 & 0 & 0 & 4 \\
\end{tabular}
\end{center}
I costi ridotti sono tutti non negativi: la soluzione ottima è $x^* = (2, 0)$, con valore $-4$. Nella tabella finale $B^{-1}$ si legge nelle colonne di $x_3$, $x_4$ e $x_5$, e la base è $(x_1, x_4, x_5)$, con $c_B = (-2, 0, 0)$:
\[
\bar{y}^T = c_B^T B^{-1} = (-2, 0, 0) \begin{pmatrix} 1 & 0 & 0 \\ -2 & 1 & 0 \\ -1 & 0 & 1 \end{pmatrix} = (-2, 0, 0).
\]
Allo stesso risultato si arriva prendendo l'opposto dei costi ridotti delle variabili di scarto, che sono $(2, 0, 0)$. Il punto $\bar{y}$ è ammissibile per il duale e ha valore $2 \cdot (-2) = -4$, uguale all'ottimo del primale\footnote{La soluzione primale è degenere, perché $x_4$ è in base con valore nullo. In questo caso il duale può avere più soluzioni ottime: anche $y = (0, -1, 0)$ è ammissibile e ha valore $-4$. Il simplesso ne trova una, che dipende dalla base finale.}.

\section{Analisi post-ottimale}\label{sec:post-ottimale}
Una volta risolto un problema, ci si può chiedere come cambia la soluzione ottima se cambiano i dati, ad esempio se aumenta la disponibilità di una risorsa. Le variabili duali rispondono proprio a questa domanda nel caso dei termini noti.

\paragraph{Esempio.}
Un'azienda produce due prodotti, deluxe e standard, utilizzando due ingredienti $I_1$ e $I_2$ e una lavorazione su macchina. Le quantità necessarie per una unità di prodotto e i prezzi di vendita sono riportati nella tabella:
\begin{center}
\begin{tabular}{lccc}
\toprule
 & Deluxe & Standard & Disponibilità settimanale \\
\midrule
$I_1$ (kg) & 3 & 2 & 1200 \\
$I_2$ (kg) & 4 & 1 & 1000 \\
Macchina (ore) & 2 & 1 & 700 \\
Prezzo (euro) & 24 & 14 & \\
\bottomrule
\end{tabular}
\end{center}
Indicando con $x_1$ e $x_2$ le quantità prodotte di deluxe e standard, il problema e il suo duale sono:
\[
\begin{aligned}
&\max\, (24x_1 + 14x_2) \\
&3x_1 + 2x_2 \le 1200 \\
&4x_1 + x_2 \le 1000 \\
&2x_1 + x_2 \le 700 \\
&x_1,\, x_2 \ge 0
\end{aligned}
\qquad\Longrightarrow\qquad
\begin{aligned}
&\min\, (1200y_1 + 1000y_2 + 700y_3) \\
&3y_1 + 4y_2 + 2y_3 \ge 24 \\
&2y_1 + y_2 + y_3 \ge 14 \\
&y_1,\, y_2,\, y_3 \ge 0
\end{aligned}
\]
ovvero lo stesso problema del primo esempio della sezione~\ref{sec:regole-duale}, con i termini noti moltiplicati per $100$. La soluzione ottima primale è $x^* = (160, 360)$, con un profitto di $8880$ euro.

Calcoliamo la soluzione duale con le condizioni di complementarietà. I primi due vincoli sono attivi in $x^*$, mentre il terzo vale $2 \cdot 160 + 360 = 680 < 700$: avanzano 20 ore di lavorazione, quindi $y_3 = 0$. Inoltre $x_1$ e $x_2$ sono positivi, quindi i due vincoli duali sono attivi:
\[
\begin{cases}
3y_1 + 4y_2 = 24 \\
2y_1 + y_2 = 14
\end{cases}
\quad\Longrightarrow\quad y_1 = 6.4, \quad y_2 = 1.2.
\]
Il duale vale $1200 \cdot 6.4 + 1000 \cdot 1.2 = 8880$, come il primale: per il corollario~\ref{cor:ottimalita-duale} le due soluzioni sono ottime.

Supponiamo ora che l'azienda voglia aumentare la disponibilità di $I_1$ o di $I_2$ per produrre e guadagnare di più. Di quanto aumenta il profitto se la disponibilità di $I_1$ passa da 1200 a 1201 kg a settimana? La regione ammissibile del duale non dipende dai termini noti del primale: questi compaiono solo nella funzione obiettivo duale, che diventa $1201y_1 + 1000y_2 + 700y_3$. Se la soluzione duale ottima resta la stessa, il valore ottimo aumenta quindi di $y_1^* = 6.4$ euro.

Allo stesso modo un kg in più di $I_2$ vale $1.2$ euro, mentre un'ora di macchina in più non vale niente, perché ne avanzano già 20.

In altre parole, la variabile duale $y_i^*$ esprime di quanto varia il valore ottimo del primale se si incrementa di una unità la risorsa $i$-esima\footnote{Nell'esempio questo vale finché la disponibilità di $I_1$ resta tra 750 e 1250 kg. Oltre questi valori la base ottima cambia e bisogna risolvere di nuovo il problema.}.

Vediamo il caso generale. Consideriamo il primale in forma standard con base ottima $B$, e perturbiamo i termini noti, che passano da $b$ a $b + \Delta b$. I costi ridotti $c_N^T - c_B^T B^{-1} N$ non dipendono da $b$, quindi restano non negativi e la soluzione duale $\bar{y}^T = c_B^T B^{-1}$ resta ammissibile. La nuova soluzione di base primale è invece $(B^{-1}(b + \Delta b), 0)$, e si presentano due casi:
\begin{itemize}
    \item \textbf{Base ancora ammissibile}: se $B^{-1}(b + \Delta b) \ge 0$, la soluzione di base primale è ancora ammissibile, quindi la base $B$ resta ottima. La soluzione duale non cambia, e il valore ottimo varia di
    \[
    c_B^T B^{-1}(b + \Delta b) - c_B^T B^{-1} b = \bar{y}^T \Delta b.
    \]
    In particolare, se si perturba solo il termine noto $b_r$, il valore ottimo varia di $\bar{y}_r \Delta b_r$.
    \item \textbf{Base non più ammissibile}: se $B^{-1}(b + \Delta b)$ ha componenti negative, la soluzione di base primale non è più ammissibile, ma quella duale sì. Si può allora ripartire dalla base $B$ con il simplesso duale (sezione~\ref{sec:simplesso-duale}).
\end{itemize}
Se $\bar{y}_r$ è grande, piccole variazioni di $b_r$ producono grandi variazioni del valore ottimo. Se invece $\bar{y}_r$ è piccolo, il valore ottimo è poco sensibile al termine noto $b_r$.

\section{Il simplesso duale}\label{sec:simplesso-duale}
Il \textbf{simplesso duale} è una variante del simplesso visto nel capitolo~\ref{ch:simplesso}, che d'ora in poi chiamiamo \textbf{simplesso primale}. Serve a risolvere problemi in cui si ha a disposizione una base che supera il test di ottimalità, ma che non è ammissibile.

\paragraph{Esempio.}
Consideriamo il problema
\[
\begin{aligned}
&\min\, (3x_1 + x_2 + x_3) \\
&2x_1 + x_2 \ge 5 \\
&2x_1 + x_2 + 3x_3 \ge 2 \\
&x_1,\, x_2,\, x_3 \ge 0
\end{aligned}
\qquad\Longrightarrow\qquad
\begin{aligned}
&2x_1 + x_2 - x_4 = 5 \\
&2x_1 + x_2 + 3x_3 - x_5 = 2 \\
&x_i \ge 0, \quad i = 1, \dots, 5.
\end{aligned}
\]
Le variabili di surplus hanno colonne con il segno meno, quindi il sistema non è in forma canonica. Si potrebbe pensare di usare il metodo delle due fasi, e l'idea non sarebbe sbagliata, ma richiede un problema artificiale con due variabili in più. Proviamo invece a moltiplicare entrambi i vincoli per $-1$:
\[
-2x_1 - x_2 + x_4 = -5, \qquad -2x_1 - x_2 - 3x_3 + x_5 = -2.
\]
Ora le colonne di $x_4$ e $x_5$ formano la matrice identità, e la tabella è:
\begin{center}
\begin{tabular}{c|ccccc|c}
 & $x_1$ & $x_2$ & $x_3$ & $x_4$ & $x_5$ & $b$ \\
\hline
$x_4$ & $-2$ & $\boxed{-1}$ & 0 & 1 & 0 & $-5$ \\
$x_5$ & $-2$ & $-1$ & $-3$ & 0 & 1 & $-2$ \\
\hline
 & 3 & 1 & 1 & 0 & 0 & 0 \\
\end{tabular}
\end{center}
La soluzione di base $(0, 0, 0, -5, -2)$ non è ammissibile, perché i termini noti sono negativi. I costi ridotti però sono tutti non negativi, quindi il test di ottimalità è superato. È la situazione opposta a quella del simplesso primale, che parte da una base ammissibile e cerca di raggiungere l'ottimalità.

Il nome deriva dall'osservazione della sezione~\ref{sec:teoremi-dualita}: una base con costi ridotti non negativi dà una soluzione duale $\bar{y}^T = c_B^T B^{-1}$ ammissibile. Il simplesso duale genera quindi una sequenza di basi con soluzione duale ammissibile e soluzione primale non ammissibile, finché non trova una base ammissibile anche per il primale. A quel punto la base supera entrambi i test ed è ottima. Il metodo si basa sull'ammissibilità duale, ed è per questo che si chiama simplesso duale.

Geometricamente, il simplesso primale si sposta tra vertici della regione ammissibile. Il simplesso duale invece parte da un punto che non appartiene alla regione ammissibile (nell'esempio l'origine, che non soddisfa i vincoli) e si avvicina al vertice ottimo attraverso soluzioni di base non ammissibili.

\subsection{Criteri del simplesso duale}
Consideriamo una tabella in forma canonica con costi ridotti $\bar{c}_j \ge 0$ per ogni $j$. Le regole del simplesso duale sono le seguenti:
\begin{itemize}
    \item \textbf{Criterio di ottimalità}: se tutti i termini noti sono non negativi, $\beta_i \ge 0$ per ogni $i$, la base $B$ è ottima. La soluzione ottima primale è $\bar{x} = (B^{-1}b, 0)$ e quella duale è $\bar{y}^T = c_B^T B^{-1}$.
    \item \textbf{Criterio di inammissibilità}: se esiste una riga $r$ con $\beta_r < 0$ e tutti gli elementi $\alpha_{rj} \ge 0$, il primale non ha soluzioni ammissibili (e il duale è illimitato). La riga $r$ si legge infatti $x_r + \sum_j \alpha_{rj} x_j = \beta_r$: con $x \ge 0$ il primo membro è non negativo, e non può essere uguale a $\beta_r < 0$.
    \item \textbf{Criterio di uscita}: esce dalla base la variabile $x_r$ con il termine noto più negativo, ovvero $\beta_r = \min\{\beta_i : \beta_i < 0\}$. In caso di pareggio si sceglie quella con indice minore.
    \item \textbf{Criterio di entrata}: entra in base la variabile $x_s$ che realizza il rapporto minimo tra i costi ridotti e i valori assoluti degli elementi negativi della riga $r$:
    \[
    \frac{\bar{c}_s}{|\alpha_{rs}|} = \min_j \left\{ \frac{\bar{c}_j}{|\alpha_{rj}|} : \alpha_{rj} < 0 \right\}.
    \]
\end{itemize}
Il pivot è l'elemento $\alpha_{rs}$, che è negativo, e le operazioni sulle righe sono le stesse del simplesso primale. Dopo il pivot il termine noto della riga $r$ diventa $\beta_r / \alpha_{rs} > 0$.

Il criterio di entrata serve a mantenere i costi ridotti non negativi, ovvero l'ammissibilità duale. Dopo il pivot i nuovi costi ridotti valgono:
\[
\bar{c}'_j = \bar{c}_j - \alpha_{rj} \frac{\bar{c}_s}{\alpha_{rs}}.
\]
Poiché $\alpha_{rs} < 0$, se $\alpha_{rj} \ge 0$ il costo ridotto non diminuisce. Se invece $\alpha_{rj} < 0$, la condizione $\bar{c}'_j \ge 0$ equivale a $\frac{\bar{c}_j}{|\alpha_{rj}|} \ge \frac{\bar{c}_s}{|\alpha_{rs}|}$, ed è garantita proprio dalla scelta del rapporto minimo.

Mettendo insieme i criteri, i passi dell'algoritmo sono:
\begin{enumerate}
    \item Si parte da una base con costi ridotti non negativi, ovvero con soluzione duale $\bar{y}^T = c_B^T B^{-1}$ ammissibile.
    \item Se tutti i termini noti sono $\ge 0$, la soluzione corrente è ottima e l'algoritmo termina. Altrimenti si sceglie la riga $r$ della variabile uscente con il criterio di uscita.
    \item Se tutti gli elementi della riga $r$ sono $\ge 0$, il primale è inammissibile e l'algoritmo termina. Altrimenti si sceglie la variabile entrante $x_s$ con il criterio di entrata.
    \item Si esegue il pivot sull'elemento $\alpha_{rs}$ e si torna al passo 2.
\end{enumerate}

\paragraph{Esempio.}
Torniamo al nostro esempio. Il termine noto più negativo è $-5$, quindi esce $x_4$. Nella sua riga gli elementi negativi sono quelli di $x_1$ e di $x_2$, e i rapporti sono:
\[
\min \left\{ \frac{3}{|-2|},\ \frac{1}{|-1|} \right\} = \min \left\{ \frac32,\ 1 \right\} = 1
\]
quindi entra $x_2$, e il pivot vale $-1$. Con le operazioni $R_1 \to -R_1$, $R_2 \to R_2 + R_1$ e $R_3 \to R_3 - R_1$, dove $R_1$ indica la riga del pivot già aggiornata, si ottiene:
\begin{center}
\begin{tabular}{c|ccccc|c}
 & $x_1$ & $x_2$ & $x_3$ & $x_4$ & $x_5$ & $b$ \\
\hline
$x_2$ & 2 & 1 & 0 & $-1$ & 0 & 5 \\
$x_5$ & 0 & 0 & $-3$ & $-1$ & 1 & 3 \\
\hline
 & 1 & 0 & 1 & 1 & 0 & $-5$ \\
\end{tabular}
\end{center}
I termini noti sono ora non negativi, e i costi ridotti sono rimasti non negativi: la soluzione $x^* = (0, 5, 0)$ è ottima, con valore $5$, ed è bastata una sola iterazione.

\paragraph{Esempio.}
Risolviamo con il simplesso duale il problema
\[
\begin{aligned}
&\min\, (12x_1 + 11x_2 + 16x_3) \\
&{-3x_1} + x_2 + x_3 \ge 3 \\
&4x_1 + x_2 + 2x_3 \ge 4 \\
&x_1,\, x_2,\, x_3 \ge 0.
\end{aligned}
\]
Sottraiamo le variabili di surplus $x_4$ e $x_5$ e moltiplichiamo entrambi i vincoli per $-1$, ottenendo $3x_1 - x_2 - x_3 + x_4 = -3$ e $-4x_1 - x_2 - 2x_3 + x_5 = -4$. Il termine noto più negativo è $-4$, quindi esce $x_5$. I rapporti sono $\frac{12}{4} = 3$, $\frac{11}{1} = 11$ e $\frac{16}{2} = 8$, quindi entra $x_1$:
\begin{center}
\begin{tabular}{c|ccccc|c}
 & $x_1$ & $x_2$ & $x_3$ & $x_4$ & $x_5$ & $b$ \\
\hline
$x_4$ & 3 & $-1$ & $-1$ & 1 & 0 & $-3$ \\
$x_5$ & $\boxed{-4}$ & $-1$ & $-2$ & 0 & 1 & $-4$ \\
\hline
 & 12 & 11 & 16 & 0 & 0 & 0 \\
\end{tabular}
\end{center}
Con le operazioni $R_2 \to -\frac14 R_2$, $R_1 \to R_1 - 3R_2$ e $R_3 \to R_3 - 12R_2$ si ottiene:
\begin{center}
\begin{tabular}{c|ccccc|c}
 & $x_1$ & $x_2$ & $x_3$ & $x_4$ & $x_5$ & $b$ \\
\hline
$x_4$ & 0 & $-\frac74$ & $\boxed{-\tfrac52}$ & 1 & $\frac34$ & $-6$ \\
$x_1$ & 1 & $\frac14$ & $\frac12$ & 0 & $-\frac14$ & 1 \\
\hline
 & 0 & 8 & 10 & 0 & 3 & $-12$ \\
\end{tabular}
\end{center}
Ora esce $x_4$, con termine noto $-6$. I rapporti sono $8 / \frac74 = \frac{32}{7}$ e $10 / \frac52 = 4$, quindi entra $x_3$. Con le operazioni $R_1 \to -\frac25 R_1$, $R_2 \to R_2 - \frac12 R_1$ e $R_3 \to R_3 - 10R_1$ si ottiene:
\begin{center}
\begin{tabular}{c|ccccc|c}
 & $x_1$ & $x_2$ & $x_3$ & $x_4$ & $x_5$ & $b$ \\
\hline
$x_3$ & 0 & $\frac{7}{10}$ & 1 & $-\frac25$ & $-\frac{3}{10}$ & $\frac{12}{5}$ \\
$x_1$ & 1 & $\boxed{-\tfrac{1}{10}}$ & 0 & $\frac15$ & $-\frac{1}{10}$ & $-\frac15$ \\
\hline
 & 0 & 1 & 0 & 4 & 6 & $-36$ \\
\end{tabular}
\end{center}
Resta negativo il termine noto di $x_1$, che quindi esce. I rapporti sono $1 / \frac{1}{10} = 10$ per $x_2$ e $6 / \frac{1}{10} = 60$ per $x_5$, quindi entra $x_2$. Con le operazioni $R_2 \to -10R_2$, $R_1 \to R_1 - \frac{7}{10} R_2$ e $R_3 \to R_3 - R_2$ si ottiene:
\begin{center}
\begin{tabular}{c|ccccc|c}
 & $x_1$ & $x_2$ & $x_3$ & $x_4$ & $x_5$ & $b$ \\
\hline
$x_3$ & 7 & 0 & 1 & 1 & $-1$ & 1 \\
$x_2$ & $-10$ & 1 & 0 & $-2$ & 1 & 2 \\
\hline
 & 10 & 0 & 0 & 6 & 5 & $-38$ \\
\end{tabular}
\end{center}
I termini noti sono non negativi, quindi la soluzione $x^* = (0, 2, 1)$ è ottima, con valore $38$. Notiamo che il valore $z_0$ (l'opposto dell'elemento in basso a destra) è passato da $0$ a $12$, $36$ e infine $38$: nel simplesso duale cresce a ogni iterazione, perché è il valore del duale, che stiamo massimizzando.

Il duale del problema è:
\[
\begin{aligned}
&\max\, (3y_1 + 4y_2) \\
&{-3y_1} + 4y_2 \le 12 \\
&y_1 + y_2 \le 11 \\
&y_1 + 2y_2 \le 16 \\
&y_1,\, y_2 \ge 0.
\end{aligned}
\]
Per un problema con soli vincoli di tipo $\ge$, la soluzione ottima duale è il vettore dei costi ridotti delle variabili di surplus nella tabella finale\footnote{Moltiplicando le righe per $-1$, le colonne delle variabili di surplus diventano colonne dell'identità, e il loro costo ridotto vale $-\bar{y}'_i$, dove $\bar{y}'$ è la soluzione duale del sistema moltiplicato per $-1$. Le variabili duali del problema di partenza sono $\bar{y} = -\bar{y}'$, quindi coincidono con i costi ridotti.}, ovvero $\bar{y} = (6, 5)$. Infatti $\bar{y}$ è ammissibile e il suo valore è $3 \cdot 6 + 4 \cdot 5 = 38$, come quello del primale. Ricordando la sezione~\ref{sec:duale-tabella}, per un problema con soli vincoli di tipo $\le$ la soluzione duale è invece l'opposto dei costi ridotti delle variabili di scarto.

\begin{osservazione}[Simplesso duale o metodo delle due fasi?]
Consideriamo un problema di minimo con vincoli di tipo $\le$ e $\ge$, di cui almeno uno $\ge$. Moltiplicando per $-1$ i vincoli $\ge$, le variabili di scarto e di surplus formano una base, e i costi ridotti iniziali coincidono con i coefficienti della funzione obiettivo:
\begin{itemize}
    \item Se i coefficienti della funzione obiettivo sono tutti non negativi, la base iniziale supera il test di ottimalità e si può usare il simplesso duale (oppure il metodo delle due fasi).
    \item Se qualche coefficiente è negativo, la base iniziale non supera il test di ottimalità, quindi il simplesso duale non si può applicare direttamente e si usa il metodo delle due fasi.
\end{itemize}
\end{osservazione}

\section*{Riferimenti}
Il materiale visto in questo capitolo comprende:
\begin{itemize}
    \item Materiale visto a lezione (30 ottobre, 6, 13 e 16 novembre 2026).
    \item Dispensa del corso \emph{Dualità}, sezioni 1--7 (definizione, tabella delle trasformazioni, interpretazione economica, teoremi fondamentali, calcolo della soluzione duale, condizioni di complementarietà, simplesso duale e analisi post-ottimale).
\end{itemize}
